% Auxiliary integral Chow calculation; no task result is registered by this file.
\paragraph{Auxiliary lemma (proved here): the codimension-three exceptional boundary in dimension seven.}
Let $k$ be an algebraically closed field of arbitrary characteristic,
let $W_1,W_2$ be smooth cubic surfaces over $k$, and put
$a=c_1(\mathcal O_{W_1}(1))$ and
$b=c_1(\mathcal O_{W_2}(1))$. Let $C/k$ be a smooth connected
projective curve, and let $K$ be any line bundle on $C$. On
\[
 S=W_1\times W_2\times C,\qquad V=\mathbb P^1\times S,
 \qquad L=\mathcal O_{\mathbb P^1}(1)
            \boxtimes\mathcal O_{W_1}(1)
            \boxtimes\mathcal O_{W_2}(1)\boxtimes K,
\]
multiplication by $c_1(L)$ is injective integrally:
\[
 c_1(L):\mathrm{CH}^2(V)\longrightarrow\mathrm{CH}^3(V).
 \tag{EC.1}
\]
For $K=K_C$, this applies to the product construction with
cycle codimension $r=3$ and local dimension $N=7$.
It keeps every torsion element of $\operatorname{Pic}(C)$; no
coefficient is reduced modulo the characteristic.

\begin{proof}
We first specify the integral Chow groups used in the calculation.
Over $k$ a smooth cubic surface $W$ is the blow-up of $\mathbb P^2$
in six points. Write $H,E_1,\ldots,E_6$ for its standard divisor
basis and $z$ for the class of a point. The projective-bundle and
codimension-two blow-up formulas give
\[
 \begin{gathered}
 \mathrm{CH}^0(W)=\mathbb Z,\qquad
 \operatorname{Pic}(W)=\mathbb ZH\oplus\bigoplus_{i=1}^6\mathbb ZE_i,
 \qquad\mathrm{CH}^2(W)=\mathbb Zz,\\
 H^2=z,\qquad E_i^2=-z,\qquad H E_i=0,
 \qquad E_iE_j=0\ (i\ne j),\\
 c_1(\mathcal O_W(1))=3H-\sum_iE_i,
 \qquad c_1(\mathcal O_W(1))^2=3z.
 \end{gathered}
 \tag{EC.2}
\]
These are integral statements in every characteristic.
For exact sources for the formulas we use
\citet[Theorems 2.4 and 2.8(3), with the displayed maps]{cyJiang2023}
and the integral projective-space-bundle formula
\citep[Lemma 42.36.2, Tag 02TX]{cyStacksProject}.
The all-characteristic cubic-surface blow-up description is the
input already used in C.3bi and C.3bj
\citep[Chapter 4, Proposition 2.4, p.~178]{D_HuybrechtsCubics}.

For completeness, applying each point blow-up after product with an
arbitrary smooth $T$ has centre $\{q_i\}\times T$ and exceptional
divisor $\mathbb P^1\times T$. Its formula adds exactly one copy of
$\mathrm{CH}^{q-1}(T)$ in degree $q$, represented by the external
product with $E_i$. The projective-bundle formula supplies the
$1,H,H^2$ summands before the blow-ups. The self-intersection and
projection formulas give the products in (EC.2) also with these
coefficients. Thus repeated application proves, without rationalizing,
the external-product ring isomorphisms needed below:
\[
 \mathrm{CH}^*(W_1\times W_2\times T)
  \simeq\mathrm{CH}^*(W_1)\otimes_{\mathbb Z}
        \mathrm{CH}^*(W_2)\otimes_{\mathbb Z}\mathrm{CH}^*(T).
 \tag{EC.3}
\]
This uses only the displayed product blow-ups; it is not a proposed
integral K\"unneth theorem for arbitrary projective varieties.

Put $h=c_1(\mathcal O_{\mathbb P^1}(1))$, $k_C=c_1(K)$ and
$c=a+b+k_C$. The projective-bundle formula expresses an arbitrary
$\gamma\in\mathrm{CH}^2(V)$ uniquely as
\[
 \gamma=\gamma_0+h\gamma_1,
 \qquad\gamma_0\in\mathrm{CH}^2(S),\quad
 \gamma_1\in\operatorname{Pic}(S).
\]
Since $h^2=0$, the equation $(h+c)\gamma=0$ gives
\[
 \gamma_0=-c\gamma_1,
 \qquad c^2\gamma_1=0.
 \tag{EC.4}
\]
By (EC.3), uniquely write
\[
 \gamma_1=d_1+d_2+\lambda,
 \qquad d_i\in\operatorname{Pic}(W_i),\quad
 \lambda\in\operatorname{Pic}(C).
\]
The component of $c^2\gamma_1$ belonging to
$\mathrm{CH}^3(W_1\times W_2)\boxtimes\mathrm{CH}^0(C)$
is $(a+b)^2(d_1+d_2)$. Because $a^2=3z_1$ and $b^2=3z_2$,
its two direct summands in (EC.3) give
\[
 \begin{aligned}
 3d_2+2(a\cdot d_1)b&=0
                    &&\text{in }\operatorname{Pic}(W_2),\\
 3d_1+2(b\cdot d_2)a&=0
                    &&\text{in }\operatorname{Pic}(W_1).
 \end{aligned}
 \tag{EC.5}
\]
Let $x=a\cdot d_1$ and $y=b\cdot d_2$, which are integers.
Pairing the two equations with $b$ and $a$ respectively yields
$3y+6x=0$ and $3x+6y=0$. Therefore $2x+y=x+2y=0$ in
$\mathbb Z$, so $x=y=0$. Equation (EC.5) and the torsion-free
divisor bases in (EC.2) then give $d_1=d_2=0$.

It remains to use $c^2\lambda=0$.
Products of two divisor classes pulled from the curve are zero,
so this reduces to
\[
 \bigl(3z_1\boxtimes1+2a\boxtimes b
                         +3\,1\boxtimes z_2\bigr)\boxtimes\lambda=0.
 \tag{EC.6}
\]
The two pure degree-two summands give $3\lambda=0$ in
$\operatorname{Pic}(C)$. In the mixed summand, choose one exceptional
coefficient $E_{1,i}\boxtimes E_{2,j}$ from the two divisor bases.
Its coefficient in $a\boxtimes b$ is $(-1)(-1)=1$, so this
coordinate of (EC.6) gives $2\lambda=0$. The identity $3-2=1$
implies $\lambda=0$ in the full Picard group. It follows that
$\gamma_1=0$, and then $\gamma_0=0$ by (EC.4), proving (EC.1).
\end{proof}

\noindent\emph{Scope of the consequence.}
This lemma supplies exactly the exceptional-image step of a
codimension-three, dimension-seven local cycle detector. The
geometric construction, actual ordinary family, relative Milnor
class and its nonzero differential must also be verified in the
combined higher-codimension proof. The lemma by itself is not
registered as a new local restriction theorem.

\paragraph{Auxiliary reduction (proved here): the exact Smith-lattice kernel.}
Let $s\ge1$, $m\ge0$ and
\[
 T=\mathbb P^m\times\prod_{i=1}^sW_i,
 \qquad A_q=\mathrm{CH}^q(T),\qquad
 b=h+\sum_i a_i,
\]
where the $W_i$ are smooth cubic surfaces and $a_i$ their
anticanonical classes. Let $C$ be a smooth connected projective
curve and let $K$ have nonzero degree $e$ on $C$. Put
$V=T\times C$ and $l=b+c_1(K)$.
The groups $A_q$ are finite free with the integral ring basis
obtained by repeating (EC.2)--(EC.3). The map
\[
 M_b:A_{s-1}\longrightarrow A_s,\qquad v\longmapsto bv,
 \tag{EC.7}
\]
is injective. If its nonzero integral Smith invariant factors are
$d_1,\ldots,d_t$, then the exact kernel of multiplication on the
whole Chow group is
\[
 \begin{split}
 \ker\bigl(l:\mathrm{CH}^{s}(V)\longrightarrow
                         \mathrm{CH}^{s+1}(V)\bigr)
  &=\ker\bigl(M_b\otimes\operatorname{Pic}^0(C)(k)\bigr)\\
  &\simeq\bigoplus_{j=1}^t\operatorname{Pic}^0(C)(k)[d_j].
 \end{split}
 \tag{EC.8}
\]
The first equality uses the external-product inclusion. The last
isomorphism depends on integral Smith bases. In particular no
unrestricted integral hard Lefschetz assertion is made.

\begin{proof}
We first prove the elementary rational injectivity needed to remove
the free part. On the cubic Chow ring over $\mathbb Q$, multiplication
by $a$ has the three-element chain
\[
 1\longmapsto a\longmapsto 3z\longmapsto0
\]
and kills the degree-one subspace $a^\perp$, where orthogonality
refers to the intersection pairing. Define a lowering operator $F$
by $F(a)=2$, $F(3z)=2a$, $F(1)=0$ and $F(a^\perp)=0$.
Give this chain squared norms $1,2,4$ and choose any positive
rational inner product on $a^\perp$, orthogonal to the chain.
Then $F$ is adjoint to multiplication by $a$, and their commutator
$[a,F]$ acts in degree $q$ by $2q-2$.

On the projective-space factor, multiplication by $h$ has lowering
operator
\[
 F(h^j)=j(m-j+1)h^{j-1}.
\]
Choose squared norms recursively by this positive coefficient.
The same adjoint identity holds and $[h,F]$ acts in degree $j$ by
$2j-m$. Tensor the inner products and add the operators over all
factors. Writing $E$ for multiplication by $b$, we obtain
$F=E^*$ and
\[
 [E,F]\big|_{A_q\otimes\mathbb R}
       =(2q-(2s+m))\operatorname{id}.
 \tag{EC.9}
\]
For a vector $v$ of weight $w$ this gives
\[
 \|Ev\|^2=\|Fv\|^2-w\|v\|^2.
 \tag{EC.10}
\]
Therefore $E$ is injective in every negative weight. Also $E^2$
is injective in every weight $w\le-2$. If $w<-2$, apply the
preceding assertion twice. If $w=-2$ and $E^2v=0$, put $u=Ev$.
The weight-zero identity (EC.10) gives $Fu=0$, and then
$\|u\|^2=\langle Ev,u\rangle=\langle v,Fu\rangle=0$.
Thus $u=0$ and then $v=0$.
In degree $q=s-1$ the weight is $-m-2$.
It follows that both $b$ and $b^2$ are injective on
$A_{s-1}\otimes\mathbb Q$, and hence on the free group $A_{s-1}$.
This argument takes place in the explicitly identical integral
Chow ring (EC.2) in every characteristic; it does not transfer an
unproved positive-characteristic cohomological assertion.

Choose $o\in C(k)$ and write $c_1(K)=e[o]+\kappa_0$, with
$\kappa_0\in\operatorname{Pic}^0(C)(k)$. The integral product
decomposition gives, uniquely,
\[
 \gamma=\alpha+q\boxtimes[o]+\eta,
 \quad
 \alpha\in A_s,\quad q\in A_{s-1},\quad
 \eta\in A_{s-1}\otimes\operatorname{Pic}^0(C)(k).
\]
If $l\gamma=0$, its component with no curve divisor and the
degree of its curve-divisor component give
\[
 b\alpha=0,\qquad e\alpha+bq=0.
 \tag{EC.11}
\]
Applying $b$ to the second equality yields $b^2q=0$; hence $q=0$.
Since $e\ne0$ and $A_s$ is free, we then have $\alpha=0$.
The remaining equation is exactly $b\eta=0$.
Conversely every such $\eta$ is killed by $l$, because products of
two divisor classes on the curve vanish. This proves the first
equality in (EC.8). Applying integral unimodular Smith changes of
basis to (EC.7) gives the last equality, retaining the full abstract
Picard group and all its primary torsion.
\end{proof}

\noindent\emph{Finite exact lattice checks.}
The retained script \texttt{check\_cubic\_exceptional\_lattices.py}
computes exact multiplication matrices in the full basis (EC.2).
It first uses integral row and column operations with unit pivots,
then finds the Smith form of the residual integer matrix. The complete sparse
matrices, residual matrices, hashes and stdout are retained.
For $r=s+1$ its nine computed cases are:
\[
\begin{array}{c|c|r|r|l}
 r&m&\operatorname{rk}A_{s-1}&\operatorname{rk}A_s
                  &\text{nonunit Smith factors}\\\hline
 2&0&1&7&\text{none}\\
 2&1&1&8&\text{none}\\
 3&0&14&51&3\\
 3&1&15&65&\text{none}\\
 3&2&15&66&\text{none}\\
 4&0&150&385&3\ (15\text{ times}),\ 9,\ 27\\
 4&1&171&535&3,\ 9\\
 4&2&172&556&3\\
 4&3&172&557&\text{none}
\end{array}
 \tag{EC.12}
\]
The run begun at 15:03:32.933566 UTC on 19 November 2026 passed
31 finite checks in 0.723674853 seconds. The earlier six-case run
and its distinct runtime are retained separately. These finite
computations concern exactly the listed lattices. For example,
(EC.8) and the $r=4,m=1$ row bound the whole exceptional kernel
by exponent nine; they do not assert that this kernel is zero
over fields of characteristic different from three, since their
Picard groups still have three-primary torsion. A separate
order-$p$ cycle argument can use the coprimality when $p\ne3$.
